\section{Results}
\label{sec:results}

This section reports the outcomes of the evaluation described in Section~\ref{sec:experiments}. The recognition results appear first, followed by the reproduction of the per-character attention visualization of Gurav~\etal~\cite{gurav2025beyondmem} on the trained model and a critical assessment of whether the Darcet~\etal~\cite{darcet2024registers} claim transfers to the CTC HTR setting. The architecture ablation, fine-tuning study, cross-dataset generalization, synthetic-data domain gap, and reproducibility experiments follow in the same order.

\subsection{Main Recognition Results}
\label{sec:res_main}

Table~\ref{tab:register_sweep} reports the register sweep on IAM at the best-validation-CER epoch. Table~\ref{tab:arch_comparison} reports the best number per architecture family under the same evaluation protocol.

\begin{table}[h]
\centering
\footnotesize
\caption{ViT-RGTS register sweep on IAM (Aachen split). Rows marked SWA use the stochastic-weight-averaged model.}
\label{tab:register_sweep}
\begin{tabular}{rrcccc}
\toprule
Registers & SWA & Best Ep. & Val CER (\%) & Test CER (\%) & Test WER (\%) \\
\midrule
0  & --         & 142 & 3.96 & 5.74 & 18.87 \\
0  & \checkmark & 124 & 4.11 & 5.90 & 19.12 \\
2  & --         & 147 & 4.13 & 5.98 & 19.39 \\
2  & \checkmark & 112 & 4.05 & 6.16 & 20.04 \\
4  & --         & 142 & 3.88 & \textbf{5.55} & \textbf{18.26} \\
4  & \checkmark & 138 & 3.97 & 5.77 & 19.05 \\
8  & --         & 129 & 4.00 & 5.94 & 19.32 \\
8  & \checkmark & 115 & 4.19 & 5.96 & 19.51 \\
16 & --         & 146 & 4.08 & 5.88 & 19.13 \\
16 & \checkmark & 120 & 4.12 & 5.88 & 19.25 \\
\bottomrule
\end{tabular}
\end{table}

\begin{table}[h]
\centering
\small
\caption{Best result per architecture family. CNN--RNN and ViT-RGTS numbers are directly comparable within each dataset; the pretrained lanes have roughly $10\times$ more parameters.}
\label{tab:arch_comparison}
\begin{tabular}{lrrcccc}
\toprule
Architecture & Params & Reg. & SWA & Best Ep. & Test CER (\%) & Test WER (\%) \\
\midrule
CNN--RNN (IAM)                 & 7.4M  & --  & \checkmark & 145 & \textbf{4.62}  & \textbf{15.70} \\
ViT-RGTS (IAM)                 & 9.5M  & 4   & --         & 142 & 5.55           & 18.26 \\
CNN--RNN (READ2016)            & 7.4M  & --  & --         & 126 & 4.06           & 18.29 \\
ViT-RGTS (READ2016)            & 9.5M  & 0   & --         & 138 & 4.73           & 20.65 \\
HTR-VT external (READ2016)     & 8.8M  & 4   & --         & 76  & 4.82           & 20.58 \\
TorchVision ViT-B/16 (IAM)     & 91.9M & 0   & --         & 149 & 15.09          & 38.19 \\
TrOCR-base (IAM)               & 90.4M & --  & \checkmark & 149 & 12.10          & 32.31 \\
\bottomrule
\end{tabular}
\end{table}

Two observations organize the rest of Section~\ref{sec:results}. First, the register-count effect on CER is small. The range across $R \in \{0, 2, 4, 8, 16\}$ is $0.24\%$ without SWA (the $R = 4$ row is slightly outside this range at $5.55\%$ due to a longer effective training window and is discussed separately) and $0.10\%$ with SWA. The three-seed reproducibility experiment in Section~\ref{sec:res_seeds} places the noise floor at $\pm 0.08\%$, so any register-count effect on CER at this dataset size is near the resolution limit of the experiment. Second, the compact CNN--RNN baseline retains the best absolute CER on IAM at $4.62\%$, and the ViT-RGTS model with $R = 4$ registers is $0.93\%$ behind at $1.29\times$ the parameter count. Recognition is not the register mechanism's contribution; the contribution is on the attention side, which the next two subsections examine in detail.

Figure~\ref{fig:cer_convergence} shows the CER convergence of the ViT-RGTS register sweep. All register counts follow the same trajectory after epoch~$25$, and the CER-versus-$R$ separation is smaller than the epoch-to-epoch fluctuation of any single run. This is the same statement Table~\ref{tab:register_sweep} makes numerically.

\begin{figure}[h]
\centering
\includegraphics[width=0.85\linewidth]{training_cer_curve}
\caption{CER convergence of the ViT-RGTS register sweep on IAM.}
\label{fig:cer_convergence}
\end{figure}

Before the attention comparisons across register counts are presented, one vocabulary point needs setting. The CTC output at each timestep is a distribution over the character alphabet together with a blank symbol, and greedy decoding emits a character exactly at the timestep where the non-blank posterior peaks. Figure~\ref{fig:ctc_posterior} illustrates this on the IAM word ``talks.'' for the ViT-RGTS ($R = 4$) model. The top panel shows the input image with red vertical lines at each character's \emph{CTC peak column} $t_c$, defined as the timestep at which the greedy decoder emits that character. The bottom panel shows the per-character posterior over the $128$ available timesteps; only a handful of columns carry non-blank probability mass, and the remaining columns are dominated by the blank symbol. Every per-character attention overlay that follows in this section is read at these $t_c$ columns.

\begin{figure}[h]
\centering
\includegraphics[width=0.85\linewidth]{fig_ctc_posterior}
\caption{CTC posterior for ViT-RGTS ($R = 4$) on the IAM word ``talks.''. Red vertical lines mark the CTC peak columns $t_c$ used to query per-character attention overlays in the following figures.}
\label{fig:ctc_posterior}
\end{figure}

\subsection{Reproducing the Per-Character Attention Visualization}
\label{sec:res_bm}

Following the visualization of Gurav~\etal~\cite{gurav2025beyondmem} Fig.~5, for each character $c$ the attention row from the corresponding patch-token query at the final Transformer layer is upsampled to the input resolution and overlaid on the input word crop with the \emph{inferno} colormap. Figure~\ref{fig:bm_booklet} shows the reproduction on the word ``booklet''. The top panel is the input word crop, and the six lower panels are the per-character attention overlays for the recognized letters \texttt{b}, \texttt{o}, \texttt{k}, \texttt{l}, \texttt{e}, \texttt{t} at their CTC peak columns. The peak shifts left-to-right along the word as the character index advances, and each map places substantial mass on the ink strokes of the character being decoded. The qualitative behaviour Gurav~\etal\ report in their Fig.~5 is reproduced end-to-end on the trained model.

\begin{figure}[h]
\centering

\includegraphics[width=0.5\linewidth]{bm_booklet_word}\\[0.2em]

\includegraphics[width=0.155\linewidth]{bm_booklet_char0_b}\hspace{0.1em}%
\includegraphics[width=0.155\linewidth]{bm_booklet_char1_o}\hspace{0.1em}%
\includegraphics[width=0.155\linewidth]{bm_booklet_char3_k}\hspace{0.1em}%

\includegraphics[width=0.155\linewidth]{bm_booklet_char4_l}\hspace{0.1em}%

\includegraphics[width=0.155\linewidth]{bm_booklet_char5_e}\hspace{0.1em}%

\includegraphics[width=0.155\linewidth]{bm_booklet_char6_t}
\caption{Reproduction of the per-character attention visualization of Gurav~\etal~\cite{gurav2025beyondmem} Fig.~5 on the ViT-RGTS model ($R = 4$). Top: input word crop (``booklet'', IAM). Bottom: per-character self-attention overlay for the recognized letters \texttt{b}, \texttt{o}, \texttt{k}, \texttt{l}, \texttt{e}, \texttt{t} at their CTC peak columns; the peak shifts left-to-right along the word. Gurav~\etal's maps come from cross-attention in a diffusion U-Net, and ours from CTC self-attention, so exact pixel alignment is not architecturally guaranteed, but the visual grammar transfers.}
\label{fig:bm_booklet}
\end{figure}

Two caveats already noted in Section~\ref{sec:rw_interp} are visible in the figure. The peak is centered near, but not exactly on, the character stroke, and each overlay retains a spatially arbitrary secondary lobe. Section~\ref{sec:res_claim} quantifies these effects and how much of the gap the register count closes.

\subsection{Does the Register-Token Claim Hold in the CTC HTR Setting?}
\label{sec:res_claim}

Darcet~\etal~\cite{darcet2024registers} make a two-part claim for classification and self-supervised ViTs: register tokens leave downstream accuracy unchanged, and register tokens improve attention interpretability by absorbing the sink pattern that would otherwise contaminate patch attention. This subsection asks whether both halves transfer when the loss changes from a per-image classification objective to a per-frame CTC objective.

\subsubsection{Recognition half of the claim: transfers cleanly}

The recognition half transfers without qualification. On IAM the CER range across $R \in \{0, \ldots, 16\}$ is $0.24\%$ without SWA and $0.10\%$ with the SWA-averaged model (Table~\ref{tab:register_sweep}). The seed-noise floor of $\pm 0.08\%$ (Section~\ref{sec:res_seeds}) puts the observed range at most three times the noise floor, so it cannot be resolved as a monotonic register effect. On READ2016 (Section~\ref{sec:res_read2016}) the range tightens further to $0.09\%$, corroborating the IAM finding on a different script and vocabulary. This is exactly the behaviour Darcet~\etal\ report for classification ViTs, and it carries over.

\subsubsection{Attention half of the claim: transfers partially}

Figure~\ref{fig:register_effect} shows the register-count effect on the attention maps produced by the model. Rows correspond to $R = 0, 4, 8, 16$; columns correspond to the six characters of the input word ``talks.''; each panel is the self-attention overlay at that character's CTC peak column, rendered with a shared colour scale across the entire grid so that visual dynamic ranges are directly comparable.

\begin{figure}[h]
\centering
\includegraphics[width=0.92\linewidth]{fig_register_effect_talks}
\caption{Register-count effect on character-level self-attention (ViT-RGTS on IAM ``talks.''). Rows: $R = 0, 4, 8, 16$. Columns: input crop, then per-character overlay at each CTC peak column, rendered with a shared colour scale across the grid. At $R = 0$ the maps for \texttt{t}, \texttt{a}, \texttt{k} show broad diffuse bands over regions far from the queried character, matching the sink pattern of Darcet~\etal\ here surfacing on CTC patch tokens. As $R$ grows, the diffusion is progressively suppressed; by $R = 8$ the peaks concentrate onto narrow vertical stripes aligned with the ink columns of the queried character.}
\label{fig:register_effect}
\end{figure}

Read row by row, the figure directly illustrates the mechanism Darcet~\etal\ describe. At $R = 0$ the query for the first character (\texttt{t}) attends broadly across the image; whole vertical bands of unrelated regions light up. At $R = 4$ the broad bands are visibly reduced and the peak concentrates on the queried character's stroke, though a secondary lobe on the same patch column persists across characters. At $R = 8$ the maps become essentially clean narrow peaks that shift with character index; at $R = 16$ the same behaviour continues with slightly more high-frequency noise. The three attention-side metrics of Section~\ref{sec:attn_eval} agree with the visual reading. Character-attention entropy on the fixed evaluation sample decreases from a median of $5.7$~bits at $R = 0$ to $5.0$~bits at $R = 8$. The spatial ordering of the peaks tightens: the correlation between character position and the location of the attention peak (measured by Spearman's $\rho$) rises from about $0.59$ at $R = 0$ to about $0.65$ at $R = 8$. The number of self-attention heads that follow left-to-right reading order also grows with $R$, from about $8$ heads at $R = 0$ to about $10$ heads at $R = 16$ (out of $8$ per layer, averaged across layers).

Where the claim does not fully transfer is the sink-absorption prediction. In a classification ViT, register tokens end up carrying more L2 norm than patches, because they take over the global-information-storage role that sinks used to serve. In the CTC ViT here, register tokens carry lower L2 norm than patches at every $R > 0$: patches remain the higher-norm carriers. Figure~\ref{fig:register_effect} shows the same asymmetry qualitatively; the residual secondary lobe on the same patch columns across every character is the surviving sink.

\subsubsection{Why the transfer is only partial}

Two structural differences explain the partial transfer. A classification ViT has a dedicated CLS query that is architecturally distinct from the patch tokens. When registers are added they take over the global-information role and the CLS-side reads clean patches. A CTC ViT has no CLS query, so the queries that read the attention are the patch tokens themselves. If a patch has developed high norm it is simultaneously a query and a key of the sink pattern, and registers cannot dissolve that self-consistency from one side alone. In addition, the CTC loss does not directly penalize an attention map that looks noisy provided the greedy decoding is correct, so there is no gradient signal driving the sinks to fully migrate onto the registers.

A data-scale reading is also worth stating. Vision Transformers are known to be data-hungry, and the two benchmarks used here (IAM at roughly $6.5$K training lines and READ2016 at roughly $8.3$K) are small by modern ViT standards. It is entirely possible that at substantially larger training scale, of the order used for TrOCR pretraining, the register mechanism reaches the full sink-absorption regime Darcet~\etal\ report for ImageNet-scale classification ViTs. The current results are consistent with the mechanism being a real intervention that is under-driven by data at HTR-benchmark scale.

The recognition half of the Darcet~\etal\ claim therefore transfers cleanly on two datasets, and the attention half transfers partially. Attention becomes measurably less diffuse and better ordered, and more heads align to reading order, but a residual sink pattern intrinsic to CTC self-attention survives. Fully suppressing it would require either an additional query-side mechanism (a CLS-style aggregate query or a small attention-entropy regularizer) or substantially more training data than is available at HTR-benchmark scale.

\subsection{Architecture Ablation}
\label{sec:res_abl_arch}

Table~\ref{tab:ablations} isolates the effect of individual design choices in ViT-RGTS. Two rows stand out.

\begin{table}[h]
\centering
\footnotesize
\caption{Architecture ablations on IAM (ViT-RGTS baseline: $R = 4$, no SWA). Each row changes one variable from the baseline.}
\label{tab:ablations}
\begin{tabular}{lrcccc}
\toprule
Variant & Params & Epochs & Val CER (\%) & Test CER (\%) & Test WER (\%) \\
\midrule
Baseline ($R = 4$)      & 9.5M  & 150 & 3.88 & 5.55 & 18.26 \\
Depth $8$               & 11.1M & 150 & 4.07 & 5.79 & 19.00 \\
Depth $4$               & 7.9M  & 150 & 3.91 & 5.93 & 19.30 \\
Dim $128$               & 5.5M  & 150 & 4.07 & 5.93 & 19.25 \\
Dim $512$               & 36.7M & 80  & 4.42 & 6.46 & 20.68 \\
Aggressive augmentation & 9.5M  & 80  & 4.58 & 6.50 & 20.89 \\
$1$-layer BiLSTM        & 6.3M  & 150 & 4.52 & 6.80 & 22.46 \\
RNN head only           & 9.4M  & 80  & 4.87 & 6.81 & 21.95 \\
CNN head only           & 5.2M  & 150 & 6.67 & 9.47 & 30.76 \\
\textbf{No CNN stem}    & 9.2M  & 80  & \textbf{72.49} & \textbf{72.44} & \textbf{125.66} \\
Legacy row-major layout & 9.2M  & 150 & 73.05 & 73.65 & 127.20 \\
\bottomrule
\end{tabular}
\end{table}

Removing the CNN stem and the legacy row-major layout both collapse the CER by more than an order of magnitude, and they fail through the same mechanism. Without the CNN stem, the encoder falls back to a raw $16 \times 16$ patch grid flattened row-major, which interleaves vertical position into the token sequence. CTC's monotonic-alignment assumption requires strictly left-to-right token order, and this row-major layout violates it: the CTC forward-backward algorithm has no valid alignment between the ground-truth character sequence and the emitted token order. Two independent architectural configurations converging on the same failure isolate this as a causal property of CTC token-formation rather than an incidental hyperparameter effect. The remaining rows are quantitative refinements. The dual head is a real contributor ($+1.26\%$ when only the RNN branch is kept, $+3.92\%$ when only the CNN branch is kept). Widening the embedding to $512$ hurts because the learning rate was tuned for width $256$, and aggressive augmentation hurts because dataset diversity is already the bottleneck at IAM scale.

\subsection{Fine-Tuning Pretrained Models}
\label{sec:res_finetune}

Table~\ref{tab:finetune} sweeps the fine-tuning strategy for the two pretrained lanes. The naive uniform-LR fine-tune on ViT-B/16 is the cleanest available demonstration of why differential-LR strategies exist for pretrained backbones. The training loss plateaus near $119$ from epoch~$3$ onward, the CER never drops below $75\%$, and the pretrained ImageNet weights are corrupted by the freshly initialized CTC head's large gradients propagated at $\eta = 10^{-3}$ across the entire backbone. Two-group differential LR cuts CER to $15.09\%$, a five-fold error-rate improvement purely from grouping the learning rate. Layer-wise LR decay produces the best TrOCR result at $12.10\%$ CER with SWA, and gradual unfreezing fails for both models.

\begin{table}[h]
\centering
\small
\caption{Fine-tuning strategy sweep on IAM. ``Naive'' uses uniform $\eta = 10^{-3}$ across backbone and head; ``Frozen'' trains only the CTC head.}
\label{tab:finetune}
\begin{tabular}{llcc}
\toprule
Model & Strategy & Test CER (\%) & Test WER (\%) \\
\midrule
ViT-B/16 & Naive uniform LR       & 73.14 & 123.71 \\
ViT-B/16 & Frozen + SWA           & 18.34 & 43.72  \\
ViT-B/16 & 2-group LR             & \textbf{15.09} & \textbf{38.19} \\
ViT-B/16 & Gradual unfreeze       & 66.20 & 93.30  \\
TrOCR    & Frozen encoder         & 71.98 & 99.10  \\
TrOCR    & LLRD                   & 12.38 & 32.60  \\
TrOCR    & LLRD + SWA             & \textbf{12.10} & \textbf{32.31} \\
TrOCR    & Gradual unfreeze       & 13.69 & 35.98  \\
\bottomrule
\end{tabular}
\end{table}

Gradual unfreezing fails for different reasons on the two models. On ViT-B/16 the large domain gap from ImageNet keeps backbone layers frozen while the head trains, so the head overfits to features the backbone will later revise; recovering from this moving-target problem within the epoch budget is difficult. On TrOCR, whose encoder is already pretrained on handwriting, the small domain gap makes the initial frozen features close to the target, and the effect is milder. Both outcomes are consistent with the general principle that gradual unfreezing helps only when the initial frozen features are close to the target domain.

\subsection{READ2016 Generalization}
\label{sec:res_read2016}

Table~\ref{tab:read2016} reruns the register sweep on READ2016. The register-count effect on CER stays within $0.09\%$, again close to the noise floor. CNN--RNN retains the best absolute CER at $4.06\%$. The HTR-VT external comparator lands at $4.82\%$ CER, matching the $R = 4$ ViT-RGTS row exactly, which independently confirms that both CTC ViT designs sit in the same performance neighbourhood on this dataset.

\begin{table}[h]
\centering
\small
\caption{READ2016 generalization. The register-count CER range is $0.09\%$, corroborating the IAM finding on a different script and vocabulary.}
\label{tab:read2016}
\begin{tabular}{lccc}
\toprule
Model & Registers & Test CER (\%) & Test WER (\%) \\
\midrule
CNN--RNN               & --  & \textbf{4.06} & \textbf{18.29} \\
ViT-RGTS               & 0   & 4.73          & 20.65 \\
ViT-RGTS               & 4   & 4.82          & 20.93 \\
ViT-RGTS               & 8   & 4.75          & 20.65 \\
ViT-RGTS               & 16  & 4.78          & 20.77 \\
HTR-VT external        & 4   & 4.82          & 20.58 \\
\bottomrule
\end{tabular}
\end{table}

\subsection{Synthetic Pretraining and the Render-to-Real Gap}
\label{sec:res_synth}

Models trained exclusively on the ${\sim}950$K synthetic corpus reach $0.5$ to $0.8\%$ CER on synthetic validation but $30$ to $46\%$ CER on real IAM test. This is an unmistakable render-to-real domain gap. Synthetic renderings have no ink bleed, no writer-specific idiosyncratic letterforms, no paper texture, and no scanning artifacts, so models trained only on them fit rendering-specific regularities that do not transfer. ViT-RGTS transfers slightly better than CNN--RNN from synthetic to real ($30.31\%$ versus $35.45\%$ at $R = 0$), plausibly because self-attention is less able to overfit low-level rendering textures than a fixed CNN filter bank. TrOCR, whose encoder was originally pretrained on real handwriting, fails catastrophically when retrained on synthetic-only data ($75.83\%$ CER), a case of pretrained-representation forgetting driven by out-of-distribution retraining. All synthetic runs were terminated by the cluster walltime before their intended epoch budget, so the numbers are lower bounds on convergence.

\subsection{Reproducibility Across Seeds}
\label{sec:res_seeds}

The seed-noise floor was measured by re-running the baseline ViT-RGTS configuration ($R = 4$, no SWA) at three different seeds. Test CER: $5.96\%$, $5.77\%$, $5.84\%$; mean $\pm$ standard deviation is $5.86 \pm 0.08\%$. Every quantitative claim in this section is interpreted against this floor.
